% GENERATED FILE. Edit canonical lessons, then run npm run generate:books.
% canonical-source-hash: fnv1a64:6a47999d9713bb7d
% canonical-lessons: LA-C111-malo, LA-C111-cupio, LA-C111-opto, LA-C111-desidero, LA-C111-voluntas

\chapter{To prefer, To desire, To wish for, To miss, Wish}
\label{ch:la-to-prefer-to-desire-to-wish-for-to-miss-wish}
\hlchaptermodality{111}

\begin{chapteropening}
\textbf{By the end of this chapter:} \emph{I can say to prefer, to desire, to wish for, to miss and a wish.}

Complete the last lesson of chapter 111: I can say to prefer, to desire, to wish for, to miss and a wish.
\end{chapteropening}

\section[mālō, mālle]{mālō, mālle — to prefer}

\label{lesson:LA-C111-malo}



\begin{quote}
Before the new one: say the Latin for it is allowed, then the Latin for to not want.
\end{quote}

\subsection*{You'll want to know: mālō, mālle}
\textbf{mālō, mālle} — \textquotedblleft{}to prefer\textquotedblright{}.

\begin{cousinweb}
\textbf{mālō} is contracted from \emph{magis}, \textquotedblleft{}more\textquotedblright{}, and \emph{volō}, \textquotedblleft{}to want\textquotedblright{} (Lewis \& Short): to want more, so to prefer.
\end{cousinweb}

\begin{grammarlens}[title={What you've built}]
One more word for this chapter.
\end{grammarlens}

\subsection*{Your turn}
\begin{itemize}
\raggedright
  \item \emph{Say it:} \emph{mālō}
  \item \emph{Say it:} \emph{mālō}, once more
  \item \emph{Say it:} say \emph{mālō}
  \item \emph{Recall:} say the Latin for it is allowed, then the Latin for to not want, then say \emph{mālō} again
\end{itemize}

\begin{tcolorbox}[breakable,colback=teal!4,colframe=teal!35!black,title={Before you move on}]
What does mālō, mālle mean? (\textquotedblleft{}To prefer\textquotedblright{}.) Say it once more.
\end{tcolorbox}
\section[cupiō, cupere]{cupiō, cupere — to desire, to long for}

\label{lesson:LA-C111-cupio}



\begin{quote}
Before the new one: say the Latin for to not want, then the Latin for to prefer.
\end{quote}

\subsection*{You'll want to know: cupiō, cupere}
\textbf{cupiō, cupere} — \textquotedblleft{}to desire, to long for\textquotedblright{}.

\begin{cousinweb}
\textbf{cupiō} is probably inherited; Lewis \& Short compare Sanskrit \emph{kup-}, \textquotedblleft{}to be stirred, to be angry\textquotedblright{}, a link still often made. English \emph{covet} comes from its family through French.
\end{cousinweb}

\begin{grammarlens}[title={What you've built}]
One more word for this chapter.
\end{grammarlens}

\subsection*{Your turn}
\begin{itemize}
\raggedright
  \item \emph{Say it:} \emph{cupiō}
  \item \emph{Say it:} \emph{cupiō}, once more
  \item \emph{Say it:} say \emph{cupiō}
  \item \emph{Recall:} say the Latin for to not want, then the Latin for to prefer, then say \emph{cupiō} again
\end{itemize}

\begin{tcolorbox}[breakable,colback=teal!4,colframe=teal!35!black,title={Before you move on}]
What does cupiō, cupere mean? (\textquotedblleft{}To desire, to long for\textquotedblright{}.) Say it once more.
\end{tcolorbox}
\section[optō, optāre]{optō, optāre — to wish for, to choose}

\label{lesson:LA-C111-opto}



\begin{quote}
Before the new one: say the Latin for to prefer, then the Latin for to desire.
\end{quote}

\subsection*{You'll want to know: optō, optāre}
\textbf{optō, optāre} — \textquotedblleft{}to wish for, to choose\textquotedblright{}.

\begin{cousinweb}
\textbf{optō} has no sure origin; the link with a Greek root for seeing that Lewis \& Short give is uncertain. English \emph{option} and \emph{adopt} come from it.
\end{cousinweb}

\begin{grammarlens}[title={What you've built}]
One more word for this chapter.
\end{grammarlens}

\subsection*{Your turn}
\begin{itemize}
\raggedright
  \item \emph{Say it:} \emph{optō}
  \item \emph{Say it:} \emph{optō}, once more
  \item \emph{Say it:} say \emph{optō}
  \item \emph{Recall:} say the Latin for to prefer, then the Latin for to desire, then say \emph{optō} again
\end{itemize}

\begin{tcolorbox}[breakable,colback=teal!4,colframe=teal!35!black,title={Before you move on}]
What does optō, optāre mean? (\textquotedblleft{}To wish for, to choose\textquotedblright{}.) Say it once more.
\end{tcolorbox}
\section[dēsīderō, dēsīderāre]{dēsīderō, dēsīderāre — to miss, to long for}

\label{lesson:LA-C111-desidero}



\begin{quote}
Before the new one: say the Latin for to desire, then the Latin for to wish for.
\end{quote}

\subsection*{You'll want to know: dēsīderō, dēsīderāre}
\textbf{dēsīderō, dēsīderāre} — \textquotedblleft{}to miss, to long for\textquotedblright{}.

\begin{cousinweb}
\textbf{dēsīderō} has no sure origin; Lewis \& Short call it doubtful. A common account builds it, like \emph{cōnsīderō}, \textquotedblleft{}to examine closely\textquotedblright{}, from \emph{dē-}, \textquotedblleft{}away from\textquotedblright{}, and \emph{sīdus}, \emph{sīderis}, \textquotedblleft{}star\textquotedblright{}, but the story behind that is unclear.
\end{cousinweb}

\begin{grammarlens}[title={What you've built}]
One more word for this chapter.
\end{grammarlens}

\subsection*{Your turn}
\begin{itemize}
\raggedright
  \item \emph{Say it:} \emph{dēsīderō}
  \item \emph{Say it:} \emph{dēsīderō}, once more
  \item \emph{Say it:} say \emph{dēsīderō}
  \item \emph{Recall:} say the Latin for to desire, then the Latin for to wish for, then say \emph{dēsīderō} again
\end{itemize}

\begin{tcolorbox}[breakable,colback=teal!4,colframe=teal!35!black,title={Before you move on}]
What does dēsīderō, dēsīderāre mean? (\textquotedblleft{}To miss, to long for\textquotedblright{}.) Say it once more.
\end{tcolorbox}
\section[voluntās, voluntātis]{voluntās, voluntātis — a wish, the will}

\label{lesson:LA-C111-voluntas}



\begin{quote}
Before the new one: say the Latin for to wish for, then the Latin for to miss.
\end{quote}

\subsection*{You'll want to know: voluntās, voluntātis}
\textbf{voluntās, voluntātis} — \textquotedblleft{}a wish, the will\textquotedblright{}.

\begin{cousinweb}
\textbf{voluntās} is built on \emph{volō}, \emph{velle}, \textquotedblleft{}to want\textquotedblright{} (Lewis \& Short), with \emph{-tās}, an ending for a quality. \emph{volō} is usually traced to Proto-Indo-European \emph{*welh1-}, \textquotedblleft{}to choose, to wish\textquotedblright{}.
\end{cousinweb}

\begin{grammarlens}[title={What you've built}]
One more word for this chapter.
\end{grammarlens}

\subsection*{Your turn}
\begin{itemize}
\raggedright
  \item \emph{Say it:} \emph{voluntās}
  \item \emph{Say it:} \emph{voluntās}, once more
  \item \emph{Say it:} say \emph{voluntās}
  \item \emph{Recall:} say the Latin for to wish for, then the Latin for to miss, then say \emph{voluntās} again
\end{itemize}

\begin{tcolorbox}[breakable,colback=teal!4,colframe=teal!35!black,title={Before you move on}]
What does voluntās, voluntātis mean? (\textquotedblleft{}A wish, the will\textquotedblright{}.) Say it once more.
\end{tcolorbox}
